\section{AI 在安全测试中的应用与局限}

\subsection{AI 改善安全的方式}

\begin{itemize}
    \item ``Shift left'' 安全比以往更易实现
    \item LLM 可以在工作流中发现问题
    \item 自动化渗透测试（automated penetration testing）
\end{itemize}

\subsection{当前的严峻局限}

\begin{warningbox}{AI SAST 的假阳性率问题}
在 AI SAST 中，假阳性率极高：Claude Code/Codex 的假阳性率可达 50--100\%（取决于漏洞类型），而传统 SAST 技术的假阳性率也在 50\%+ 以上。这意味着大量的安全告警是误报，会严重干扰开发者的工作流。
\end{warningbox}

\begin{itemize}
    \item \textbf{基准不现实}：现有 benchmark 往往不反映真实场景，难以评估 LLM 的安全分析能力
    \item \textbf{非确定性分析}：同一 prompt 运行多次得到不同结果——如何确保捕获所有漏洞？
    \item \textbf{Context rot}：并非所有 context 都同等有效
    \item \textbf{Compaction 信息丢失}：为适应 context window 而进行的摘要压缩可能丢失关键安全细节
\end{itemize}

\subsection{开放问题}

\begin{enumerate}
    \item 如何降低漏洞检测中的假阳性和 hallucination？
    \item 如何验证 LLM 生成的补丁是安全的且不引入回归？
    \item LLM 能否解释\textbf{为什么}标记某个漏洞或提出某个修复方案？
    \item 什么是衡量 LLM AppSec 性能的正确 benchmark？
    \item LLM 应该如何嵌入 CI/CD 流程而不会用噪音淹没团队？
    \item 如果 AI 生成的补丁引入了漏洞，谁应该负责？
\end{enumerate}

\subsection{把 AI 安全检查嵌入开发流程}

课程内容落到工程实践时，最关键的不是再加一个“智能扫描器”，而是把它放到合适的位置：

\begin{table}[H]
\centering
\begin{tabular}{p{0.18\textwidth} p{0.24\textwidth} p{0.34\textwidth}}
\toprule
阶段 & 适合的 AI 安全动作 & 主要目标 \\
\midrule
编码阶段 & IDE 内联提示、局部静态检查 & 尽早阻止明显的危险模式进入代码库 \\
提交阶段 & PR 级漏洞复核、依赖变更扫描 & 在合并前拦截高风险改动 \\
测试阶段 & 自动化 DAST / fuzzing 辅助 & 发现运行时路径和环境相关问题 \\
上线后 & 监控告警摘要、事故复盘辅助 & 缩短发现与定位时间，形成经验闭环 \\
\bottomrule
\end{tabular}
\caption{AI 安全能力在开发流程中的嵌入点}
\end{table}

\begin{knowledgebox}{最有效的策略通常不是 ``让 AI 独立负责安全''}
更现实的做法是把 AI 放在“高召回、低决策权”的位置：它负责尽可能早地暴露风险，人类和硬性测试系统负责最终裁决。这样既能利用 AI 的覆盖面，又能控制假阳性和幻觉带来的噪音成本。
\end{knowledgebox}

\subsection{本章小结}

AI 在安全测试领域既是机遇也是挑战。LLM 可以降低``shift left''的门槛，但假阳性率高、非确定性输出和 context 管理问题仍需解决。责任归属是一个尚未回答的关键伦理问题。

